%
% Discussion.tex
%
% Copyright (C) 2025 by Inspirefly.
%
% 2.305 GHz Circular Patch Antenna Documentation
%

\chapter{Discussion}

\section{Interpretation of Results}

The analytical cavity model predicts the required radius to within 0.5~mm of the HFSS-tuned value, which confirms that Eq.~\ref{eq:fres} with the fringing correction (Eq.~\ref{eq:ae}) is adequate for initial sizing at 2.305~GHz on FR-408HR. The early $\lambda_g/2$ square-patch estimate gave a correct order of magnitude but does not describe the circular $\mathrm{TM}_{11}$ resonance; its retention in git is useful only as a record of the initial bounding exercise.

The MATLAB particle-swarm run recorded in \texttt{best.txt} did not produce a viable antenna. The combination of high weights on gain and axial ratio with a single-layer, lossless substrate model drove the optimizer toward a small truncated-corner patch with a feed near the geometric centre, where the $\mathrm{TM}_{11}$ mode has a current minimum and input resistance near zero. The resulting RL${}=0$~dB and $-6.97$~dBi gain are therefore not a failure of the concept but a consequence of an under-constrained search space and an idealized feed model. This outcome motivated the shift to direct parametric sweeps in HFSS, where the feed radius $\rho_0$ and patch radius $a$ are varied independently and the full multilayer stack is represented.

HFSS results show that a single orthogonal feed pair, when driven in quadrature externally, achieves the two requirements that the MATLAB model could not satisfy simultaneously: 50~$\Omega$ match and axial ratio below 3~dB. The hexagonal patch variant confirmed that reducing symmetry introduces mode splitting; the additional $\mathrm{TM}_{21}$ resonance near 3.8~GHz coupled into the operating band and increased cross-polarization, which is why that branch was not carried into KiCAD.

\section{Future}

The PCB was fabricated on 30~August~2026 at the final 1.3414~mm FR-408HR stack (Table~\ref{tab:stackup}); what remains uncertain are the effects of the fab. tolerances (refer to JLCPCB), and ensuing data measurement. The team is considering the HackRF Pro as a low-cost VNA: its 1~MHz to 6~GHz SDR with VNA firmware/software can directly measure $S_{11}$ and inter-port $S_{21}$ for the 50~$\Omega$ match, and the same hardware can then act as transceiver for over-the-air S-band link tests in both linear (one port terminated) and circular (quadrature via the external divider) configurations, with divider/cable loss de-embedded in the VNA calibration.

If the center frequency is offset, adjusting radius $a$ in $0.1$~mm steps in HFSS and a DXF re-export will restore centre frequency without touching the feed network. Longer-term, parent-bus integration and thermal-vacuum testing will verify that the thin prepreg maintains adhesion and that the U.FL connectors survive vibration.

\nomenclature{VNA}{Vector Network Analyzer}
\nomenclature{SDR}{Software-Defined Radio}